\section{Step 1: Purpose \& Requirements}

\subsection{Purpose of the IoT System}
The purpose of the Smart Mosque Management System is to design an intelligent IoT-based solution that enhances the operational efficiency of mosques by:
\begin{itemize}
  \item Optimizing energy consumption (lighting and air conditioning)
  \item Conserving water usage in wudu (ablution) areas
  \item Monitoring and managing crowd occupancy
  \item Improving indoor air quality and comfort
\end{itemize}

The system leverages real-time sensor data, automated control mechanisms, and a centralized monitoring dashboard to ensure efficient resource utilization, especially during high-attendance periods such as Friday prayers and Ramadan.

\subsection{Functional Requirements}
These define what the system must do.

\subsubsection*{A. Energy Management}
\begin{itemize}
  \item Automatically turn ON/OFF lights and fans/AC based on:
  \begin{itemize}
    \item Prayer times (via API)
    \item Occupancy detection (motion sensors)
  \end{itemize}
  \item Adjust energy usage dynamically depending on crowd density.
\end{itemize}

\subsubsection*{B. Occupancy Detection}
\begin{itemize}
  \item Count number of people entering/exiting using IR sensors.
  \item Provide real-time hall capacity status.
  \item Trigger alerts if capacity exceeds safe limits.
\end{itemize}

\subsubsection*{C. Water Monitoring (Wudu Area)}
\begin{itemize}
  \item Measure water flow using flow sensors.
  \item Detect abnormal usage (leaks or taps left open).
  \item Automatically close water supply using water valves.
\end{itemize}

\subsubsection*{D. Air Quality Monitoring}
\begin{itemize}
  \item Measure air quality (e.g., CO$_2$, humidity, temperature).
  \item Trigger ventilation (fans) when thresholds are exceeded.
\end{itemize}

\subsubsection*{E. Central Dashboard}
\begin{itemize}
  \item Display:
  \begin{itemize}
    \item Current occupancy
    \item Energy consumption
    \item Water usage
    \item Air quality readings
  \end{itemize}
  \item Allow manual override control for administrators.
\end{itemize}

\subsubsection*{F. Notifications \& Alerts}
\begin{itemize}
  \item Send alerts for overcrowding, water leakage, poor air quality, and system faults.
\end{itemize}

\subsection{Non-Functional Requirements}
These define how well the system should perform.

\subsubsection*{A. Performance}
\begin{itemize}
  \item Real-time response (low latency in sensor-actuator loop)
  \item Efficient handling of multiple sensors simultaneously
\end{itemize}

\subsubsection*{B. Scalability}
\begin{itemize}
  \item System should support small to large mosques
  \item Easy addition of new sensors/devices
\end{itemize}

\subsubsection*{C. Reliability}
\begin{itemize}
  \item Continuous operation with minimal downtime
  \item Fault-tolerant communication between devices
\end{itemize}

\subsubsection*{D. Security}
\begin{itemize}
  \item Encrypted communication using WPA2-PSK for wireless sensor-to-gateway connections
  \item Dashboard access control with username/password authentication
  \item MAC filtering to allow only authorized IoT nodes
\end{itemize}

\subsubsection*{E. Usability}
\begin{itemize}
  \item Simple and intuitive dashboard interface
  \item Minimal training required for mosque staff
\end{itemize}

\subsubsection*{F. Maintainability}
\begin{itemize}
  \item Easy software/hardware updates
  \item Modular design for replacing faulty components
\end{itemize}

\subsubsection*{G. Cost Efficiency}
\begin{itemize}
  \item Use affordable IoT components suitable for local deployment in Egypt
\end{itemize}

\subsection{Constraints and Assumptions}
\subsubsection*{Constraints}
\begin{itemize}
  \item Limited budget; use cost-effective sensors
  \item Implementation in Cisco Packet Tracer simulation
  \item Dependence on internet for prayer-time API
\end{itemize}

\subsubsection*{Assumptions}
\begin{itemize}
  \item Stable local network exists in mosque
  \item Power supply is reliable
  \item Admin has basic technical knowledge
\end{itemize}

The system aligns with Egypt's sustainability goals by reducing unnecessary consumption of water and electricity in high-usage public religious spaces.
